% =============================================================================
% Annexure
% =============================================================================

\appendix
\renewcommand{\thesection}{Annexure~\Alph{section}}
\renewcommand{\thesubsection}{\Alph{section}.\arabic{subsection}}

\section{Formalization Amendments}
\label{app:amendments}

We resolved 36 design questions in a pre-registered formalization document before any code was written. Empirical work surfaced five amendments, each documented with the experiment that motivated it. We summarize them here; the full formalization (including unchanged questions) is released alongside the implementation.

\paragraph{A.1 ($\delta_{\max}$ must scale with model size):} The original specification of $\delta_{\max}=0.001$ as a per-parameter clip bound saturates within $\sim 30$ training steps on a 7B model with rank-16 LoRA: 36/36 LoRA parameters hit the clip on every step thereafter, and validation always fails. Required correction: $\delta_{\max} \approx 0.01$ at 7B scale. The per-parameter bound is the wrong abstraction; what matters is the Frobenius norm of the LoRA update relative to the model's natural weight scale. A scaling rule of $\delta_{\max}(d) = 0.001 \cdot \sqrt{d / 768}$ partially corrects this but underestimates the empirically required value.

\paragraph{A.2 ($\alpha_{\mathrm{slow}}$ fails for two separable reasons):} The original specification of $\alpha_{\mathrm{slow}} = 10^{-5}$ produces gradient updates that round to zero in bfloat16 storage: training loss is completely flat across 100 steps. Bumping to $10^{-4}$ produces a phase transition (loss $1.46 \to 0.55$). We subsequently isolated whether this is purely a dtype artefact by re-running the original default in float32 on the same benchmark. It is not. In float32 the loss does move---but only from $2.902$ to $2.874$ over 100 steps ($\sim$1\%), with $0/12$ validation passes and $\mathrm{BCP} = 0.998$, i.e.\ the adapter is effectively unmoved. There are therefore \emph{two stacked failures}: a genuine bfloat16 precision floor (updates round to exact zero) \emph{and} a learning rate that is under-scaled for a 7B model in any precision. Required correction: $\alpha_{\mathrm{slow}}$ must clear both bounds---$\alpha_{\mathrm{slow}} \geq 10 \cdot \mathrm{ulp}(\mathrm{dtype})$ for representability, and a magnitude that scales with model size for effect. The original $10^{-5}$ implicitly assumed both float32 storage and a smaller model.

\paragraph{A.3 ($\mu_{\mathrm{surprise}}$ has a generator-discriminator gap):} The quality check ``replay must be at least as surprising as the model's baseline'' rejects all replays in practice, because replays are generated from the same model that defines the surprise function and are therefore systematically more in-distribution than external calibration text. The empirical surprise of generated replays ($\sim 1.25$ nats) falls below the baseline computed on calibration text ($\sim 1.41$ nats). Required correction: replace the absolute threshold with a relative-gain criterion ($\mathrm{surprise(\mathrm{replay})} > \beta \cdot \mathrm{surprise(\mathrm{seed})}$) or compute the threshold from self-generated samples. Until amended, we disable the absolute check.

\paragraph{A.4 (Validation has a data-volume floor):} The validation criterion $\mathrm{surprise}_\mathrm{new} < (1 - \epsilon_\mathrm{learn}) \cdot \mathrm{surprise}_\mathrm{old}$ with $\epsilon_\mathrm{learn} = 0.10$ is achievable only above a minimum number of replays. With 5 facts and 2--3 replays, $0\%$ of candidates pass; with 200 facts and $\sim 50$ replays, $40\%$--$60\%$ pass depending on safety settings. The criterion itself is sound, but the formalization should explicitly document that pass rates are zero below a benchmark-dependent volume threshold.

\paragraph{A.5 (Tags are pointers, not storage):} The original formalization stored only the K/V projections of tagged tokens, but tagged tokens are short novel-entity mentions (4--8 tokens) that capture the location of the surprise without its content. Required correction: when a tag fires, store K/V for the full sentence containing the tagged span (the \emph{episode}), not just the tagged tokens. This corresponds to the proposal's own framing of tags as pointers but was implemented as direct storage. Section~\ref{sec:results_episodes} reports the empirical correction.

\section{Hyperparameter Table}
\label{app:hyperparams}

\begin{table}[ht]
\centering
\caption{Full hyperparameter list. ``Source'' indicates which of the 36 formalization questions specifies the parameter. ``Status'' indicates whether the value matches the original specification or was empirically amended.}
\label{tab:hyperparams_full}
\begin{tabular}{llll}
\toprule
Parameter & Value & Source & Status \\
\midrule
Tagging $\kappa$ (z-score threshold) & 1.5 & Q1.2 & default \\
Tagging $g$ (gap tolerance) & 3 & Q1.3 & default \\
Tagging min span length & 4 & Q1.3 & default \\
Tagging EMA $\beta$ & 0.99 & Q1.2 & default \\
PRP weights $(w_E, w_a, w_x, w_r)$ & (0.35, 0.30, 0.15, 0.20) & Q2.2 & default \\
PRP budget multiplier $c_\mathrm{prp}$ & 500 & Q2.4 & default \\
PRP cross-ref similarity threshold & 0.5 & Q2.2 & default \\
LoRA rank ($r$) & 16 & Q3.1 & default \\
LoRA $\alpha$ & 32 & Q3.1 & default \\
Adapted layer fraction & 0.333 (top 1/3) & Q3.2 & default \\
Adapted matrices & V, O & Q3.1 & default \\
$\alpha_\mathrm{fast}$ & $10^{-4}$ & Q3.3 & default (LoRA W$_\mathrm{fast}$ only) \\
$\alpha_\mathrm{slow}$ & $10^{-4}$ & Q3.3 & \textbf{amended (A.2)} \\
$\delta_\mathrm{max}$ & 0.01 & Q3.4 & \textbf{amended (A.1)} \\
$\lambda_\mathrm{EWC}$ & 100 & Q3.4 & default \\
$\phi_\mathrm{min}$ (plasticity floor) & 0.10 & Q3.4 & default \\
$\epsilon_\mathrm{learn}$ (validation) & 0.10 & Q4.6 & default \\
$\epsilon_\mathrm{degrade}$ (rollback) & 0.02 & Q4.5 & default \\
Replay temperature & 0.7 & Q4.1 & default \\
Replay top-$p$ & 0.9 & Q4.1 & default \\
Compression target & 5 & Q4.1 & default \\
Steps per memory & 5 & Q4.4 & default \\
Sleep weight decay & 0.01 & Q4.4 & default \\
KV bank capacity (tokens) & 20{,}000 & --- & new (A.5) \\
KV top-$k$ retrieval gate & 16 (default), 64 (episode) & --- & new (A.5) \\
KV memory storage unit & full sentence (episode) & --- & \textbf{amended (A.5)} \\
$\mu_\mathrm{surprise}$ (quality check) & disabled (set to 0) & Q4.1 & \textbf{amended (A.3)} \\
\bottomrule
\end{tabular}
\end{table}

\section{Implementation Details}
\label{app:implementation}

\subsection{KV Memory Injection Mechanics}

KV memory injection requires care to avoid silently wrong attention scores. We document the choices made:

\textbf{Position encoding.} We treat memory entries as occupying negative position IDs $[-n_\mathrm{mem}, -1]$. RoPE~\citep{su2024roformer} is applied at read time to memory keys with these positions. The encoding is mathematically valid (sin is odd, cos is even, both well-defined at negative inputs) and verified unit-tested. Memory thus appears to current queries (at positive positions) as ``recent context tokens,'' a pattern the model was trained on through RoPE's relative-position mechanic.

\textbf{Attention mask.} Qwen2 uses float additive masks with $-\infty$ for masked positions and $0.0$ for visible. When extending the mask with memory columns, we fill those columns with $0.0$ of the mask's own dtype. A boolean True or integer 1 here would silently produce wrong scores on float-mask attention paths.

\textbf{Cache ordering.} Memory K/V must be injected \emph{after} the \texttt{past\_key\_values} cache update, not before. If memory is injected first, the cache accumulates a duplicate copy of memory tokens on every autoregressive generation step; by step 50 the cache contains 50 copies of memory and BCP collapses (we observed $\sim 96\times$ degradation under this bug).

\textbf{K/V extraction layer-norm.} When extracting K/V from a forward pass via \texttt{output\_hidden\_states=True}, the apparent ``hidden state at layer $\ell$'' is the input to layer $\ell$, but attention computes K/V from the layer-normed input. Skipping the layer-norm produces silently wrong K/V. We add an integration test that captures K/V via a forward hook during a normal forward pass and asserts equality with our extraction.

\subsection{Per-Phase KV State}

A correctness requirement we discovered: KV injection must be \emph{disabled} during W$_\mathrm{cons}$ training. If left enabled, the model can satisfy the training loss by relying on injected memory rather than encoding the replay content into weights---a crutch the model would not have at inference time after the bank is cleared. We add a \texttt{set\_enabled(bool)} method to the injector and toggle as follows:

\begin{itemize}[itemsep=0.1em,topsep=0.2em]
    \item Wake (tagging, KV write): enabled.
    \item Phase 1 Generate (replay): enabled (so the model attends to source while compressing).
    \item Phase 2 Quality check: state preserved (no parameter changes; result invariant).
    \item Phase 3 Train (W$_\mathrm{cons}$ updates): \emph{disabled} (forces W$_\mathrm{cons}$ to encode into weights).
    \item Phase 4 Perplexity check: disabled (clean preservation measurement).
    \item Phase 5 Validate \& cleanup: disabled (clean recall measurement).
    \item After successful consolidation: bank is cleared, enabled flag reset for next cycle.
\end{itemize}

\subsection{Memory Gate Placement}

The warm-up's memory gate (Section~\ref{sec:warmup_method}) is applied to the memory \emph{values} immediately before they are concatenated onto the key/value time axis, and after the memory keys have received their negative-position RoPE. Placement matters for a reason worth recording: attention output is linear in $V$ but non-linear in $K$, so gating values scales exactly the memory contribution to the residual stream, whereas gating keys would perturb the softmax allocation itself. A consequence of this choice is that the gate cannot reclaim attention mass---memory positions still absorb softmax weight even at $g \to 0$---so it modulates the strength of the memory contribution rather than switching it off. Because the gate is initialised to the identity, installing it is a provable no-op: our first test asserts bit-identical logits against the un-gated path, which caught the alternative placements during development.

\section{Multi-Cycle Configuration}
\label{app:multi_cycle}

The multi-cycle experiment of Section~\ref{sec:multi_cycle} uses Setting A from the Pareto sweep ($\delta_{\max} = 0.02$, $\lambda_{\mathrm{EWC}} = 50$, $\alpha_{\mathrm{slow}} = 10^{-4}$, steps-per-memory $= 4$) for the SLEEP arm. The naive LoRA arm uses 200 gradient steps per cycle on the cycle's batch only, with AdamW (lr $10^{-4}$, weight decay 0.01, batch size 32, uniform-with-replacement sampling). Both arms accumulate adapter state across cycles; weights from cycle $i$ persist into cycle $i+1$. Both arms are evaluated on identical control texts and recall benchmarks after each cycle. Replay strategy for the SLEEP arm is set to \texttt{original} (replays are the original fact tokens), eliminating replay quality as a confound and isolating the multi-cycle dynamics of the consolidation training step. The ten-cycle programme divides the 200 facts into 10 disjoint batches of 20 and repeats the paired protocol with three seeds on Qwen2.5-7B and two seeds on each of Qwen2.5-1.5B, Llama-3.1-8B, and Mistral-7B (nine run-pairs). The earlier three-cycle configuration used batches of $\approx 67$, and we note that batch size per cycle is a candidate moderator of the plateau onset.

\section{Validation-Matrix Protocol}
\label{app:matrix}

The revised-pipeline validation of Section~\ref{sec:results_matrix} was executed as a staged campaign whose gates, predictions, and stop-loss were committed to the repository before any run: (D0) the three-check ground-truth verification gate on all four mid-MLP configurations; (D1) per-family recipe tuning at seed 0, advancement criterion trained-subset DRA $\geq 0.10$ at $\mathrm{BCP} \leq 1.5$, plus a mid-campaign-registered clip ladder for Llama with its selection rule fixed before execution; (D2) single-cycle, five seeds per model at the tuned recipes (200 facts, 1{,}500 training steps); (D3) ten-cycle, both arms, two seeds per model (20 facts/cycle, 400 distillation steps/cycle for the SLEEP arm, matched mid-MLP adapter geometry for the naive arm); (D4) the matched-arm EWC-only comparison (online EWC, $\lambda=100$, Fisher accumulated over each cycle's facts and re-anchored after consolidation, at most 20 Fisher samples per cycle). The v2-moderate operating point throughout is $\delta_{\max}=0.1$ (0.03 for Llama per its ladder), plasticity scaling on, $\lambda_{\mathrm{EWC}}=0$ within the sleep-train step, $\alpha_{\mathrm{slow}}=10^{-4}$. The campaign comprised 35 runs on a single RTX A6000 in one 19.7-hour queue with zero failures; pre-registration documents, per-run JSON results, and execution logs are released in the repository (\texttt{logbook/} and \texttt{experiments/results/}).

\section{Reproducibility Notes}
\label{app:repro}

\paragraph{Seeding.} A single entry point seeds Python's \texttt{random}, NumPy, PyTorch (CPU and CUDA), and \texttt{PYTHONHASHSEED}. Multi-seed experiments are executed one seed per process so that model initialisation is genuinely independent rather than sequentially derived; the aggregating harness collects per-run JSON and reports mean and sample standard deviation.

\paragraph{Environment sensitivity.} Three environment issues materially affected results and are recorded so that replication does not rediscover them. First, \texttt{transformers} must be pinned to the version documented in the repository: a later release imports a distributed-tensor symbol unavailable in the CUDA image's PyTorch build, and fails at import. Second, adapter modules must be constructed in the wrapped layer's dtype; a float32 adapter attached to a bfloat16 model raises a dtype mismatch at the first matrix multiply. Third, the patched attention forward must not assume architecture-specific attributes: \texttt{sliding\_window} exists on Qwen2 attention modules but not on Llama's or (in the pinned release) Mistral's, and is read defensively. The cross-architecture verification harness of Section~\ref{sec:setup} caught the third issue before any invalid run executed; all are recorded in the released logs.

\paragraph{Verification practice.} Every modification to the KV-injection path is accompanied by a hook-based equality test against the model's own computation, written before the feature it verifies. This practice caught two silent-failure bugs during development: a missing layer-norm in K/V extraction, and a cache-ordering error that duplicated memory on every generation step. The memory gate's identity-initialisation test follows the same pattern.
